% Results section draft: B1 to B4 subsections (Phase 8, Shravan's plan section 8 item 2).
% Drafted by Dhruv for Section 4 of the QST manuscript. Satyabrat owns the structure of Section 4 and the Stage 6/7
% subsections; Shravan owns framing. This file reports data only: no interpretation, prior-work comparison or
% novelty wording. Every number is listed with its committed source in results_b1_b4_numbers.csv and checked by
% tests/test_results_draft.py. Figure labels match paper/figures (fig1 to fig6). Notation follows the reports:
% S = F_+ + F_-, r = (F_- - F_+)/S, b_X = slope of the median log10 of X against n (decades per qubit).

\subsection{The shot-ratio identity and the slope rule on entangling circuits}
\label{sec:results-b1}

\emph{Claim tested: the identity $M_{\mathrm{SWAP}}/M_{\mathrm{LE}} = (2-Q)/(S-Q)$ holds per parameter setting on
every circuit, and the exponent gap $b_{\mathrm{SWAP}} - b_{\mathrm{LE}}$ equals $b_S$ whenever $b_r \approx 0$.}

We tested both statements on 17 cells (figure~\ref{fig:b1}): the RX-product benchmark ($n = 2$--$20$) and two
entangling families, a hardware-efficient ring ansatz and an RX--RY brickwork with CZ gates, each at depth 2, 4 and
$n$, with the differentiated parameter in the early, middle or late layer ($n = 4$--$16$, or $4$--$14$ at depth $n$).
Each cell uses 10 seeds of 4,000 parameter settings per $n$; intervals are bootstrap 95\% CIs of the 10-seed mean.

\paragraph{Identity.} The identity held in all 6,740,000 per-setting rows, with maximum relative error
$8.3\times10^{-15}$ on the entangling circuits and $1.75\times10^{-11}$ on the RX product. No row was excluded.

\paragraph{Fidelity and relative-gradient slopes.} $b_S$ depends on family and depth but not on the parameter
position (positions agree to $\le 0.001$): 0.376, 0.319 and 0.303 for the ring ansatz at depth 2, 4 and $n$, and
0.382, 0.340 and 0.305 for the brickwork, against 0.601 for the RX product. The relative gradient does not
concentrate: $|b_r| \le 0.0029$ and $|b_r|/b_S \le 0.0096$ in all 16 entangling cells. Six $b_r$ intervals exclude
zero; the largest is the brickwork depth-$n$ middle cell, $-0.0029$ $[-0.0036, -0.0023]$, a slight growth of $|r|$
with $n$.

\paragraph{Gap rule.} Over the full range the measured gap is slightly below $b_S$ in every cell: the residual
$\varepsilon_{\mathrm{gap}} = (b_{\mathrm{SWAP}} - b_{\mathrm{LE}}) - b_S$ lies between $-0.0048$ and $-0.0015$
decades per qubit ($-1.28\%$ to $-0.48\%$ of $b_S$, figure~\ref{fig:b1}(b)), and its CI contains zero in none of the
16 entangling cells. For $n \ge 8$ it lies between $-0.0039$ and $+0.0005$, and the CI contains zero in 8 of 16 cells.
A post-hoc decomposition splits the residual exactly into an $O(S)$ term, which vanishes once the median $S$ is
below about $10^{-2}$, and a term from taking a difference of medians rather than a median of differences; the second
term persists at about $-0.003$ in the depth-2 and depth-4 cells over $n = 8$--$16$.

\paragraph{Exponent ratio.} The observed ratio $b_{\mathrm{SWAP}}/b_{\mathrm{LE}}$ is 1.969--2.010 over the full
range and 1.977--2.019 for $n \ge 8$, against 1.996--2.020 predicted from each cell's own $b_S$ and $b_r$
(figure~\ref{fig:b1}(d)). The per-qubit bases are circuit-dependent: about 2.4, 2.1 and 2.0 for Loschmidt and 5.7,
4.4 and 4.0 for SWAP at depth 2, 4 and $n$ (ring ansatz), not the 4 and 16 of the RX product.

\paragraph{Limits.} Two ansatz families, one target state, noiseless circuits and $n \le 16$. The fits use 6 or 7
values of $n$.

\subsection{Seed-to-seed uncertainty of the RX-product exponents}
\label{sec:results-b3}

\emph{Claim tested: the Stage~7 exponents are stable across independent parameter populations.}

We regenerated the RX-product benchmark with 20 independent seeds of 25,000 parameter settings ($n = 2$--$20$) and
report the seed mean with its 95\% percentile interval over seeds (figure~\ref{fig:landscape}(c)). For SNR $\ge 1$ the
Loschmidt slope is 0.6053 $[0.6034, 0.6070]$ (seed SD 0.0010) and the SWAP slope is 1.2014 $[1.1985, 1.2050]$
(seed SD 0.0019); every seed has $R^2 \ge 0.9998$. The Stage~7 values, 0.6059 and 1.2025, lie inside these
intervals, as do all ten fitted Stage~7 targets. Extending the fits to $n = 2$--$24$ gives slopes of
0.6002--0.6051 (Loschmidt) and 1.1985--1.2022 (SWAP) across the fitted targets.

The paired gap $b_{\mathrm{SWAP}} - b_{\mathrm{LE}}$ is 0.5961 (seed-mean CI $[0.5956, 0.5966]$) over $n = 2$--$20$,
below $\log_{10}4 = 0.6021$. For $n = 8$--$20$ it is 0.6008: the seed percentile interval $[0.5976, 0.6049]$
contains $\log_{10}4$, while the CI of the 20-seed mean, $[0.5998, 0.6017]$, lies just below it.

\subsection{Conditional sign accuracy of the Loschmidt gradient}
\label{sec:results-sign}

\emph{Claim tested: a nonzero Loschmidt gradient estimate has the correct sign with probability $(1+|r|)/2$ when
$MS \ll 1$.}

On the entangling circuits (16 cells at $n = 12$, five quantiles of $S$, $M = 1$ to 65,536; 720 exact evaluations,
nothing fitted) the deviation from $(1+|r|)/2$ is positive in every row and grows with $MS$
(figure~\ref{fig:sign}(a,b)). Its maximum is $9\times10^{-5}$ for $MS \le 10^{-3}$, $4.8\times10^{-4}$ for
$MS \le 10^{-2}$, $4.5\times10^{-3}$ for $MS \le 0.1$ and 0.045 for $MS \le 1$; above $MS = 1$ it reaches 0.33.

On the RX product the predicted population median is $(1+1/\sqrt2)/2 = 0.8536$. With cancellation-free
probabilities, the observed median at $n = 20$ is 0.8527--0.8530 for $M$ up to 16,384 (figure~\ref{fig:sign}(c)).
At $n = 20$, $M = 1024$ the seed mean is 0.8547 with seed percentile interval $[0.8494, 0.8592]$; the reference
median, 0.85267, equals the median of $(1+|r|)/2$ over the same 20,000 settings to five digits. Where $MS$ is not
small the observation departs from the prediction, as expected for a small-$MS$ statement: at $n = 12$, $M = 4096$ the
observed median is 0.868.

\subsection{Four estimators: shot exponent and fixed-budget failure}
\label{sec:results-b2}

\emph{Claim tested: the shot exponent of an estimator, and the way it fails at a fixed budget, can be measured
separately.}
% Whether this is a headline claim (Main 3) depends on P2/P4; the data below stand either way.

We compared four unbiased estimators of the same fidelity gradient on the RX-product benchmark: Loschmidt, ancilla
SWAP, destructive SWAP (two copies, no ancilla) and a Hadamard-test U-statistic. Cost is counted in total circuit
executions per gradient component ($2M$ for the first three, $4M$ for Hadamard).

\paragraph{Exponents.} For SNR $\ge 1$ over $n = 2$--$20$ (figure~\ref{fig:estimators}(a,b); 10 replicate seeds,
95\% percentile intervals) the slopes are 0.6059 for Loschmidt (replicates 0.6058 $[0.6048, 0.6064]$), 1.2025 for both
SWAP readouts (1.2022 $[1.2003, 1.2045]$) and 0.6066 for Hadamard (0.6063 $[0.6051, 0.6075]$). The two SWAP readouts
agree to $2\times10^{-13}$ in their medians. Hadamard needs 9.0 times the executions of Loschmidt at $n = 2$ and
9.4--9.55 times for $n \ge 4$ (figure~\ref{fig:estimators}(c)). An entangling spot check (ring ansatz, depth $n$,
$n = 6$--$12$) gives the same grouping: 0.302 for Loschmidt and Hadamard, 0.604 for SWAP.

\paragraph{Fixed budgets.} At $n = 14$ and $T = 32$ to 32,768 total executions (figure~\ref{fig:estimators}(d--f)),
Loschmidt returns an exactly zero estimate with median probability at least 0.9996, and when it is nonzero its sign is
correct with median probability 0.856--0.866. Both SWAP readouts rarely return zero (0.140 falling to 0.0044) but
their nonzero estimates have the correct sign with probability 0.500. Hadamard also rarely returns zero (0.158 falling
to 0.0004) and its nonzero estimates are likewise correct with probability 0.500--0.501. Figure~\ref{fig:failure}
places each estimator by its exponent and its failure statistic at $T = 2048$.

\paragraph{Limits.} One landscape and one $n$ for the fixed-budget panels. Executions are counted per gradient
component; qubit count, depth, ancillas and state copies differ between readouts, so equal executions do not mean
equal hardware cost. Noiseless simulation only.

\subsection{Depolarizing noise}
\label{sec:results-b4}

\emph{Claims tested: the identity and the slope rule hold for noisy trial states; whether the exponent doubles under
noise.}
% Whether "the noise model decides doubling" is a headline claim (Main 4) depends on P2.

Global depolarizing noise after each of $L$ layers maps $F \to qF + (1-q)/2^n$ with $q = (1-p)^L$; local single-qubit
depolarizing noise keeps the RX product a product state. Both overlays were checked against density-matrix simulation
(maximum error $6.7\times10^{-16}$ global, $5.1\times10^{-14}$ local) and reuse the B1 populations
(figure~\ref{fig:noise}).

\paragraph{Identity and slope rule.} The identity held on all 41.2 million noisy rows (maximum relative error
$1.75\times10^{-11}$, none excluded). The gap residual $\varepsilon_{\mathrm{gap}}$ stayed between $-0.012$ and
$+0.006$ decades per qubit in every cell at every $p$.

\paragraph{Global noise on the RX product.} The SWAP exponent is unchanged (1.201 over $n = 2$--$20$ at every $p$),
while the Loschmidt exponent rises: over $n \ge 8$, $b_{\mathrm{LE}}$ is 0.717, 0.803, 0.868 and 0.895 for
$p = 10^{-4}$, $10^{-3}$, $10^{-2}$ and $10^{-1}$ (figure~\ref{fig:noise}(a)). Three-point local slopes over
$n = 16$--$20$ reach $b_{\mathrm{LE}} = 0.90$ and an exponent ratio of 1.334 at $p = 10^{-1}$, against the
asymptotes $\log_{10}8 = 0.903$ and $4/3$. The crossover $n^*$, where the noise floor equals the median signal
scale, is 15.4, 12.1, 8.7 and 5.2 for $p = 10^{-4}$ to $10^{-1}$ (figure~\ref{fig:noise}(d)).

\paragraph{Local noise on the RX product.} $b_r$ stays at zero ($-0.0001$) and $b_S$ falls to 0.5834, 0.5451 and
0.4451 for $p = 10^{-3}$, $10^{-2}$ and $10^{-1}$, against 0.5831, 0.5447 and 0.4449 from the closed form
(figure~\ref{fig:noise}(c)). The exponent ratio stays at 2.00, 2.00 and 1.99.

\paragraph{Global noise on the entangling cells.} At $p = 10^{-1}$ the ratio over $n \ge 8$ falls to 1.74--1.93
across the 16 cells, with $b_r$ between 0.012 and 0.053. In the depth-$n$ cells the measured $b_r$ follows the
prediction (ring ansatz, early parameter: 0.051 against 0.048). The fixed-depth cells are mostly pre-asymptotic: the
crossover falls inside $n \le 16$ only at $p = 10^{-1}$.

\paragraph{Limits.} Simulated depolarizing channels with an ideal SWAP target copy. Local noise on entangling
circuits, other channels and hardware are not covered.
